\cotefile{158}
\scanpage{2}
% l. 1
$\unsure{\mathcal X|\mathscr X|X}$ topos loc. connexe
% ll. 2–5
\diagram{\begin{tikzcd}[ampersand replacement=\&,column sep=large]
 B \arrow[r,"{\unsure{p|P}}"] \arrow[rr,bend left=27,"\pi"] \& \mathcal X \arrow[r,"\pi_0"] \& \mathrm{Ens}
\end{tikzcd}}
\margin{stable par $\varprojlim$ de type $\Phi$}
\inserted{$\Phi$-exact à g.}

% ll. 6–7
\[
 I\longmapsto F^I\deleted{\unsure{\to\mathrm{Ens}|\in\mathrm{Ens}}}
 \qquad
\]
\[
 \operatorname{Hom}(G,F^I)=\operatorname{Hom}(I,\operatorname{Hom}(G,F))
\]
% ll. 8–10
\diagram{\begin{tikzcd}[ampersand replacement=\&,column sep=large,row sep=large]
 {\overset{\deleted{\unsure{C|G}}}{Z}} \arrow[r,"i"] \arrow[d,"p"'] \& C \arrow[dl,"q"'] \arrow[r,"{\deleted{\unsure{Q|\mathcal L}}P}"] \& \mathcal X\\
 B \arrow[ur,dashed,bend right=16,"\sigma"'] \arrow[urr,dashed,bend right=9,"{\deleted{\unsure{P|p}}\unsure{=|\simeq}P\circ\sigma}"'] \& \&
\end{tikzcd}}
$P$ est $\Phi$-exact à g. ?

% ll. 11–14
\[
 P\left(b,\underset{\substack{\rotatebox{-90}{$\in$}\\\mathcal P(\pi_0(p(b))}}{E}\right)=\widetilde E
\]
$\simeq$ ens des \unsure{sommandes|sommants} directs $\widetilde E$ de $p(b)$
\diagram{\begin{tikzcd}[ampersand replacement=\&,row sep=small]
 {} \arrow[d,dash] \& {} \& y' \arrow[d,dash,dashed]\\
 {} \& x\arrow[r] \& y
\end{tikzcd}}
% ll. 15–17
\[
\begin{aligned}
 &\deleted{\bigl(b_{i},E_{\unsure{i|1}}\bigr)}\ P(b,E)\times P(b',E')=\\
 &\hspace{2em}\coprod_{(z,z')\in E\times E'}
 \underset{\substack{\rotatebox{-90}{$\subset$}\\\pi(b)}}{\widetilde z}
 \times\underset{\substack{\rotatebox{-90}{$\subset$}\\\pi(b')}}{\widetilde {z}'}
 \subset \underbrace{\pi(b)\times\pi(b')}_{\times(E\times E')}\\
 &=\pi(b\times b')\times(E\times E')
\end{aligned}
\]
% archive sheet number: 1

\scanpage{3}
% letterhead
\begin{center}
ACADEMIE D'AIX-MARSEILLE\\
\rule{4em}{0.4pt}\\
INSPECTION ACADEMIQUE\\
DE VAUCLUSE\\
\rule{1.6em}{0.4pt}
\end{center}
% ll. 4–7, left and right
\[
 C_{c_\infty}\xrightarrow{u'_{i*}}C_{b_i}
 \qquad \deleted{\unsure{\pi|\Pi}(b)\mathbin{\unsure{-|\to}}\pi(b')}
\]
\[
 u'_{i*}\deleted{^{\unsure{*|+}}}
 \qquad E\xrightarrow{u}E'
\]
\[
 \deleted{\unsure{u'|b'}}\quad c_\infty\longrightarrow u'_i{}^{*}(c_i)
 \qquad \mathcal P^*(E)\longrightarrow\mathcal P^*(E')
\]
\[
 c\longrightarrow\qquad\qquad
 \begin{aligned}
 u(F)&\subset F'\\
 &\unsure{\Updownarrow|\parallel}\\
 F&\subset u^{-1}(F')
 \end{aligned}
\]
% ll. 8–10
\[
 \overset{\unsure{\leftarrow|\sim}}{\coprod}_{z\in\underset{\substack{\rotatebox{-90}{$\subset$}\\\unsure{\pi_b|\Pi_b|\pi(b)}}}{E}}z
 \times\coprod_{z'\in E'}z'
 =\coprod_{(z,z')\in E\times E'}\underbrace{z\times z'}_{\text{produit dans }Z^\wedge}
\]
% ll. 11–15, crossed calculations
\deleted{\unsure{est la somme des|est la somme de}}
\[
 \deleted{\unsure{=|?}}\qquad
 \deleted{\unsure{\coprod|\pi}(b\times b')\times}\quad
 \text{(\unsure{produits|produit} externes)}
\]
\[
 \deleted{\unsure{\int|d}\ \coprod\quad z\times}
 \qquad \deleted{z\times z'}\quad\text{\unsure{non|…}}
\]
\diagram{\begin{tikzcd}[ampersand replacement=\&]
 \deleted{z} \arrow[d,dash] \& \deleted{z'}\arrow[d,dash] \& \deleted{z\times z}\\
 \deleted{b}\& \deleted{b'}\&
\end{tikzcd}}
\diagram{\begin{tikzcd}[ampersand replacement=\&,row sep=small]
 \& b\times b^{\unsure{\prime|*}}\arrow[dl]\arrow[dr] \&\\
 b \& \& b
\end{tikzcd}}
% archive sheet number: 2

\scanpage{4}
% letterhead
\begin{center}
ACADEMIE D'AIX-MARSEILLE\\
\rule{4em}{0.4pt}\\
INSPECTION ACADEMIQUE\\
DE VAUCLUSE\\
\rule{1.6em}{0.4pt}
\end{center}
% l. 4
\diagram{\begin{tikzcd}[ampersand replacement=\&]
 B\arrow[r,"\pi"]\arrow[rr,bend right=24,"\widetilde\pi"'] \& \mathrm{Ens}\arrow[r,"\mathcal P^*"] \& \mathrm{Ord}
\end{tikzcd}}
% ll. 5–7
\diagram{\begin{tikzcd}[ampersand replacement=\&,row sep=small]
 C=\widetilde\pi\backslash B\arrow[d,"\text{fidèle}"]\\ B
\end{tikzcd}}
\[
\begin{gathered}
\underbrace{\operatorname{H\unsure{o|{\uline{o}}|\underset{C}{o}}m}\bigl((b,\underset{\substack{\rotatebox{-90}{$\in$}\\
 \widetilde\pi(b)=\mathcal P^*(\pi(b))}}{E}),(b',\underset{\substack{\rotatebox{-90}{$\in$}\\
 \mathcal P^*(\pi(b'))=\widetilde\pi(b')}}{E'})\bigr)}
 \longleftrightarrow\operatorname{Hom}_B(b,b')\\
 \rotatebox{90}{$=$}\\
 \{u:b\longrightarrow b'\mid \widetilde\pi(u)(\unsure{\ddot E|E})\subset E'\}
\end{gathered}
\]
% ll. 8–9
NB $\widetilde\pi\backslash B$ est la catégorie cofibrée sur $B$ définie par le foncteur \deleted{$\unsure{\mathcal P\pi|\mathcal P^*\pi}$}\inserted{composé} $: B\xrightarrow{\widetilde\pi}\mathrm{Ord}\longrightarrow\mathrm{Cat}$

% ll. 10–14
\[
 i\longmapsto c_i=(b_i,E_i)
\]
\diagram{\begin{tikzcd}[ampersand replacement=\&] I\arrow[r]\&C\arrow[r]\&B \end{tikzcd}}
\[
 b=\varprojlim b_i\longrightarrow b_i
\]
\diagram{\begin{tikzcd}[ampersand replacement=\&,column sep=large]
 c \arrow[d,dash]\arrow[r,"\text{cocart}"]\arrow[rr,bend left=17,"v_i"] \&
 {\underset{\deleted{\unsure{c_\infty=(b_i,E_i)|c_i=(b_i,E_i)}}}{u_{\infty*}(c)=c_\infty}}
 \arrow[d,dash]\arrow[r,"w_i"] \& c_i=(b_i,\unsure{\xi_i|E_i})\\
 {\deleted{\unsure{b_i|b}}b}\arrow[r,"u_\infty"]\arrow[rr,bend right=21,"u_i"'] \&
 b_\infty\overset{\text{déf}}{=}\varprojlim b_i\arrow[r,"u'_i"] \& b_i
\end{tikzcd}}
\[
 w'_{i*}(c_\infty)\xrightarrow{\text{ds }C_{b_i}}\deleted{\unsure{\xi|c}}c_i
\]
% archive sheet number: 3

\scanpage{5}
% ll. 1–4
\[
 B\xrightarrow{\pi}\mathrm{Ens}
\]
\[
 \widetilde\pi\backslash B\simeq B/\widetilde\pi=C
\]
\diagram{\begin{tikzcd}[ampersand replacement=\&,row sep=small]
 \pi\backslash B=Z\arrow[d,"\text{coétale}"]\\B
\end{tikzcd}}
\[
 C\longrightarrow\deleted{\pi\backslash\unsure{B|P}} Z^\wedge
\]
\[
 \left(b,\underset{\substack{\rotatebox{-90}{$\in$}\\\mathcal P(\deleted{\unsure{\pi|Z}}Z_b)}}{E}\right)
 \longmapsto\coprod_{z\in E}\deleted{\unsure{q|\xi}}z
\]
% ll. 5–9
\[
\begin{aligned}
 (b,E)&\longrightarrow(b',E) & u_*^{E,E'} &: E\longrightarrow E'\\
 E\subset Z_b&\longrightarrow Z_{b'}\supset E' & \deleted{\unsure{b|B}}z&\longmapsto z'\\
 b&\xrightarrow{u}b'
\end{aligned}
\]
\[
 \begin{aligned}
 \widetilde E&=\coprod_{z\in E}\overset{\unsure{\leftarrow|\sim}}{z}\deleted{\unsure{\sim|\to}}\longrightarrow\widetilde E'\\
 &&z\longrightarrow\widetilde z'
 \end{aligned}
\]
% ll. 10–12
\diagram{\begin{tikzcd}[ampersand replacement=\&,column sep=large,row sep=large]
 Z\arrow[r,hook,dash,"{i}","{\text{cocart}}" description,near start]\arrow[d,"p","\text{cofibré}"']\arrow[rr,bend left=40,"{\unsure{c_Z|\varepsilon_Z}}","{\deleted{\unsure{i|z}}}"'] \&
 {\deleted{\unsure{Z|z}}C}\arrow[r,"\varphi"]\arrow[dl,"q" description,"{\substack{\text{cofibré et}\\\text{fibré}}}"'] \&
 Z^\wedge\arrow[d,"p_!"]\\
 \uline{B}\arrow[rr,"{c_B}"'] \& \& B^\wedge
\end{tikzcd}}
% ll. 13–14
\diagram{\begin{tikzcd}[ampersand replacement=\&,row sep=small]
 \deleted{\unsure{\varphi|p}}\& b\times b'\arrow[dl,"p_1"']\arrow[dr,"p_2"]\&\\
 b\&\&b'
\end{tikzcd}}
\[
\begin{aligned}
 (b,E)\times(b',E')&=(b\times b',p_1^*(E)\cap\\
 &\hspace{7em}p_2^*(E'))
\end{aligned}
\]
% archive sheet number: 4

\scanpage{7}
\mspage{\textcircled{1}}
% ll. 1–5
$X\in\operatorname{Ob}\mathrm{Cat}$ tel que tout objet \unsure{0-connexe|0-conn.} de $X^\wedge$ soit 1-connexe

$F\in\operatorname{Ob} X^\wedge$, $U=\operatorname{Im}(F\longrightarrow e_{X^\wedge})$. Donc $F\in\operatorname{Ob}X^\wedge_{/U}\unsure{\simeq|\cong}(\underbrace{X_{/U}}_Y)^\wedge=\deleted{\unsure{X|x}}Y^\wedge$, et en \deleted{\unsure{tant|tel}} qu'objet de $Y^\wedge$, $F\longrightarrow e_{Y^\wedge}$ épi.

% ll. 6–8
\deleted{\unsure{$I$|Soit}} Soit $\widetilde\Delta={\rm Cat.}$ des ens. finis de la forme $\widetilde\Delta_n=[0,n]\cap\mathbb Z$ (équiv. à la cat. des ens finis non vides)
% ll. 9–11
\[
\begin{aligned}
 I&\longmapsto\pi_0(F^I)\ :\ \widetilde\Delta^{\mathrm o}\longrightarrow(\mathrm{Ens}),\quad\text{soit}\\
 \pi_F&=\text{\uline{$\pi\in\widetilde\Delta^\wedge$}}\ :\qquad\pi(I)=\pi_0(F^I)\qquad\text{i.e.}\\
 &&\pi(\widetilde\Delta_{\unsure{n|m}})=\pi_0(F^{n+1}\deleted{_{\unsure{*|1}}})
\end{aligned}
\]
% ll. 12–14
Je vais définir un foncteur
\[
 \boxed{\unsure{\varphi|\Phi}:(\widetilde\Delta_{/\pi})^{\mathrm o}\longrightarrow Y^\wedge=(X_{/U})^\wedge\approx X^\wedge_{/U}}
\]
% ll. 15–18
Un objet de $\widetilde\Delta_{/\pi}$ est un couple $(\widetilde\Delta_n,\alpha)$, avec $\alpha\in\pi(\widetilde\Delta_n)=\pi_0(F^{n+1})$, on pose
\[
 \varphi(\widetilde\Delta_n,\alpha)=\underbrace{\alpha}_{\substack{\text{composante}\\\text{conn. de }F^{n+1}}}\subset\deleted{\unsure{e|\alpha}}F^{n+1}
\]
\margin{Comme $F$, les $F^I$, sont ds $X^\wedge_{/U}\overset{\approx}{\longrightarrow}Y^\wedge$ de même $\alpha$.}
% ll. 19–22
Une flèche de $\widetilde\Delta_{/\pi}$ \deleted{est}, de $(\widetilde\Delta_n,\alpha)$ vers $(\widetilde\Delta_m,\beta)$, est formé d'une flèche $f:\widetilde\Delta_n\longrightarrow\widetilde\Delta_m$ telle que $\underset{\substack{\rotatebox{90}{$=$}\ {\rm def}\\\pi(\rho)(\beta)}}{{\rho}_\pi^*(\beta)}=\alpha$. Or
% archive sheet number: 6

\scanpage{8}
% Read after turning the sheet through 180 degrees.
% letterhead
\begin{center}ACADEMIE D'AIX-MARSEILLE\\\rule{4em}{0.4pt}\\
INSPECTION ACADEMIQUE\\DE VAUCLUSE\end{center}
% ll. 1–4
\deleted{$X\in\operatorname{Ob}\mathrm{Cat}$, tel que les objets 0-conn. de $X^\wedge$ soient 1-connexes.}

\deleted{$F\in\operatorname{Ob}X^\wedge$}\qquad \deleted{d'où}
\[
 \deleted{I\longmapsto F^I}\qquad
 \deleted{\left(\overbrace{\underbrace{\unsure{\mathrm{Ensf}_*|\mathrm{Ens}f_*}}_{\substack{\text{Ens finis}\\\text{non vides}}}}^{\widetilde\Delta}\right)^{\mathrm o}\longrightarrow X^\wedge}
\]
% ll. 5–7
\deleted{d'où}
\[
 \deleted{\underset{\substack{\rotatebox{90}{$=$}\\\underline{\operatorname{Hom}}(\widetilde\Delta,\mathrm{Ens})=\widetilde\Delta^\vee}}{(\widetilde\Delta^{\mathrm o})^\wedge}\longrightarrow X^\wedge}
\]
\deleted{foncteur qui commute aux $\varinjlim$}

% ll. 8–13
\deleted{Soit $U\subset e_{X^\wedge}$ l'image de $F\longrightarrow e_{X^\wedge}$, alors}
\diagram{\begin{tikzcd}[ampersand replacement=\&]
 \deleted{\widetilde\Delta^{\mathrm o}}\arrow[r]\arrow[dr]\& \deleted{X^\wedge}\\
 \&\deleted{(X_{/U})^\wedge}\arrow[u]
\end{tikzcd}}
\deleted{se factorise par}
\margin{\deleted{foncteur image directe par le \unsure{morphisme|morphism} de \deleted{\unsure{$\varphi$|$\phi$}} topos $X^\wedge_{/U}\longrightarrow X^\wedge$}\quad \deleted{$\underset{(X_{/U})^\wedge}{\rotatebox{90}{$\simeq$}}$}}
\deleted{d'où}
\[
 \deleted{\widetilde\Delta^\vee\longrightarrow(X_{/U})^\wedge\simeq X^\wedge_{/U}}
\]
\margin{\deleted{commutant aux $\varinjlim$ en plus des $\varprojlim$}}
% ll. 14–16, below the struck block
\[
 Y^\wedge\xrightarrow{f}Z^\wedge
\]
\[
 G\xrightarrow{\unsure{\Psi|\psi}\deleted{\unsure{f|\varphi}}}\unsure{G|h}
\]
\[
 f^*(G^\Psi)= f^*(G)^{f^*(\unsure{\varphi|\Psi})}
\]
% archive sheet number: 7

\scanpage{9}
\mspage{\textcircled{2}}
% letterhead
\begin{center}
\deleted{ACADEMIE D'AIX-MARSEILLE}\\
\rule{4em}{0.4pt}\\
\deleted{INSPECTION ACADEMIQUE}\\
\deleted{DE VAUCLUSE}\\
\rule{1.6em}{0.4pt}
\end{center}
% ll. 1–5
$\rho$ définit $\rho^*:F^{\widetilde\Delta_m}\longrightarrow F^{\widetilde\Delta_n}$, et la relation $\rho_\pi^{\unsure{*|\prime}}(\beta)=\alpha$ signifie que $\rho^*\mid\beta$ se factorise par $\alpha$, i.e. $\rho^*$ induit
\[
 \unsure{\varphi|\Phi}(\rho):
 \underset{\substack{\rotatebox{90}{$=$}\\\varphi(\widetilde\Delta_m,\beta)}}{\beta}\longrightarrow
 \underset{\substack{\rotatebox{90}{$=$}\\\varphi(\widetilde\Delta_n,\alpha)}}{\alpha}.
\]
% ll. 6–10
On a ainsi défini une loi fonctorielle $\varphi$. Le foncteur $\varphi$ se prolonge par continuité en
\[
 \Phi:\underbrace{\bigl(\overbrace{(\widetilde\Delta_{/\pi})^{\mathrm o}}^{Z}\bigr)^\wedge}
 _{\substack{\rotatebox{90}{$=$}\\(\widetilde\Delta_{/\pi})^\vee=Z^\wedge}}
 \longrightarrow Y^\wedge,
\]
foncteur commutant aux $\varinjlim$ quelconques
% ll. 11–15
Je suspecte que \uline{ce foncteur est exact à gauche},$^*$ donc qu'il correspond à un morphisme de topos
% slanted
\margin{$^*$ c'est ici que pourrait intervenir l'hypothèse de \unsure{connexité|connexion} faite \unsure{sur|en} $X$ (et \unsure{laquelle est|qui est-elle} \unsure{requise : ?|requise ?}), mais pas encore}
\[
 f:Y^\wedge\longrightarrow Z^\wedge
\]
\[
 (\text{où}\qquad Z\overset{\text{déf}}{=}(\widetilde\Delta_{/\pi})^{\mathrm o}).
\]
% archive sheet number: 8

\scanpage{10}
\mspage{3}
% ll. 1–9
Ce morphisme est $W_0$-asphérique, i.e. pour tout $z\in\operatorname{Ob}Z$, son image inverse $f^*(z)=\deleted{\unsure{\Phi|\varphi}}\varphi(z)$ $(=\alpha$, si $z=(\widetilde\Delta_n,\alpha))$ est 0-connexe, ce qui implique que pour tout $G\in\operatorname{Ob}Z^\wedge$ qui est 0-connexe, \deleted{$\unsure{\Phi|\varphi}$} $f^*(\deleted{\unsure{z|F}}G)$ est 0-connexe. (Pour le voir, on \inserted{\unsure{doit|devra|peut}} utiliser le fait que $f^*=\Phi$ commute aux produits fibrés, ou tout au moins aux intersections de sous-objets).
% ll. 9–15
Si on fait appel maintenant à l'hyp. que les objets 0-conn. de $X^\wedge$ (donc aussi ceux des topos induits \unsure{$-/^\wedge$|$Y^\wedge$}) sont 1-connexes, il s'ensuit que $f$ est \deleted{\unsure{un|=}} $\unsure{\mathcal W_1|W_1}$-asphérique, \deleted{i.e.} \inserted{plus précisément} que pour \unsure{tt|$\forall$} \deleted{$z\in\operatorname{Ob}Z$} $G\in\operatorname{Ob}Z^\wedge$ qui est 0-connexe, $f^*(G)$ est 1-connexe.
% ll. 16–21
Mais alors, ayant un morphisme 0-connexe $\deleted{\unsure{X|Y}}Y^\wedge_{/f^*(G)}\longrightarrow Z^\wedge_{/G}$, avec la source 1-connexe, on en conclut que le but $Z^\wedge_{/G}$ l'est aussi. Donc tout objet 0-connexe de $Z^\wedge$ est 1-connexe.
% archive sheet number: 9

\scanpage{11}
\mspage{4}
% letterhead
\begin{center}
ACADEMIE D'AIX-MARSEILLE\\
\rule{4em}{0.4pt}\\
INSPECTION ACADEMIQUE\\
DE VAUCLUSE\\
\rule{1.6em}{0.4pt}
\end{center}
% ll. 1–2
Mais je conjecture que l'on doit avoir ceci :
% ll. 3–9
\uline{Lemme} (?) Soient $\pi\in\operatorname{Ob}\widetilde\Delta^\wedge$, $Z=(\widetilde\Delta_{/\pi})^{\mathrm o}$ (donc $Z^{\mathrm o}$ une catégorie étale sur $\widetilde\Delta$ \deleted{\unsure{)|$^{\mathrm o}$}}, i.e. $Z$ coétale sur $\widetilde\Delta^{\mathrm o}$), et supposons que \inserted{a)} tout objet 0-connexe de $Z^\wedge$ soit 1-connexe et \inserted{b)} $\pi(\widetilde\Delta_{\deleted{\unsure{n|0}}0})$ est ponctuel, alors $\pi\simeq e_{\widetilde\Delta^\wedge}$ i.e. \inserted{$\pi_n=$}$\pi(\widetilde\Delta_n)$ est uniponctuel pour tout $n\in\unsure{\mathbb N|\mathcal N}$.
% ll. 10–16
Vu ce qui précède, cela implique

\uline{Corollaire 1} Soit $X$ un objet de $\mathrm{Cat}$ tel que tout objet de $X^\wedge$ qui est 0-connexe soit 1-connexe. Alors \deleted{$\unsure{X|Y}$} pour tt objet 0-connexe $F$ de $X^\wedge$, les produits cartésiens $F^n$ $(n\geq1)$ sont 0-connexes, en particulier $F\times F$ est 0-connexe.
% ll. 17–20
Vu \deleted{le fait} que \inserted{l'on sait déjà que} l'hyp. implique que $X^\wedge$ est loc. irréductible (i.e. $X$ satisfait $\unsure{\mathrm{PS}^*1|\mathrm{P5}^*1}$), on \unsure{tr.|trouve|voit} que $X^\wedge$ est loc. tot. 0-connexe, i.e. $X$ \inserted{loc.} \unsure{cofiltrante|cofiltr.te}, donc
% ll. 21–23
\uline{Corollaire 2} L'hyp. du cor 1 implique que $X$ est \inserted{loc.} \unsure{cofiltrante|cofiltr.te|cofiltrant} (et inversement, bien sûr), i.e. les comp. conn. sont \unsure{cofiltrantes|cofiltr.tes}.
% archive sheet number: 10

\scanpage{12}
\mspage{5}
% ll. 1–3
$A\in\operatorname{Ob}\mathrm{Cat}$ (je pense au cas $A=\widetilde\Delta$)

$\pi\in\operatorname{Ob}A^\wedge$, $Z=(A_{/\pi})^{\mathrm o}\simeq\pi\backslash\underbrace{A^{\mathrm o}}_B=\pi\backslash B$

$[\pi:\underbrace{A^{\mathrm o}}_B\longrightarrow\mathrm{Ens}]$ On s'intéresse à $Z^\wedge$, et à la sous-catégorie pleine de $Z^\wedge$ \deleted{\unsure{engendré|engendr.|engendrée}} stable par $\varprojlim$ finies, engendrée par $Z$. J'ai en vue un candidat : pour tout $b\in\operatorname{Ob}B$, considérons $Z_b\simeq\pi(\uline{b})$, et pour tte partie \inserted{\uline{non vide}} $E$ de $\pi(b)$
% ll. 8–12
\diagram{\begin{tikzcd}[ampersand replacement=\&,row sep=small]
 Z\arrow[d]\& Z_b\simeq\pi(b)\\ B\ni b\&
\end{tikzcd}}
\[
 E\in\mathcal P^*(\pi(b))
\]
considérons l'élément
\[
 \widetilde E=\coprod_{a\in E}a\qquad(\text{somme prise \unsure{ds|dans} }Z^\wedge)
\]
% ll. 13–16
Si $u:b\longrightarrow b'$ ds $B$, on a \inserted{$u_*=\pi(u):$} $\underset{\substack{\rotatebox{90}{$=$}\\\pi(b)}}{Z_b}\longrightarrow\underset{\substack{\rotatebox{90}{$=$}\\\pi(b')}}{Z_{b'}}$ et si $E\subset Z_b$, $E'\subset Z_{b'}$ sont tels que $u_*(E)\subset E'$, alors $u$ induit
\[
 u_*^{E,E'}:\widetilde E\longrightarrow\widetilde E'
\]
% ll. 17–19
Ainsi, le foncteur composé
\[
 \widetilde\pi:b\longmapsto\mathcal P^*(\pi(b)):\quad B\xrightarrow{\pi}\mathrm{Ens}\xrightarrow{\mathcal P^*}\mathrm{Ord}
\]
définit une catégorie $\widetilde\pi\backslash B$, \deleted{\unsure{ayant d|ayant q|dont}} $=\unsure{S|s}$
% archive sheet number: 11

\scanpage{13}
\mspage{6}
% letterhead
\begin{center}
ACADEMIE D'AIX-MARSEILLE\\
\rule{4em}{0.4pt}\\
INSPECTION ACADEMIQUE\\
DE VAUCLUSE\\
\rule{1.6em}{0.4pt}
\end{center}
% ll. 1–9
mêmes objets que $[\widetilde\pi]\backslash B$ (où $[\widetilde\pi]$ est le foncteur composé de $\widetilde\pi:B\longrightarrow\mathrm{Ord}$ avec le foncteur oubli $\mathrm{Ord}\longrightarrow\mathrm{Ens}$), mais plus de flèches, une flèche de $(b,E)$ dans $(b',E')$ $(E\in\widetilde\pi(b),E'\in\widetilde\pi(b'))$ étant une flèche $u:b\longrightarrow b'$ telle que \deleted{$\unsure{u(E)|u_*E}$} $\widetilde\pi(u)(E)\subset E'$ (\unsure{ald|ald.} $\widetilde\pi(u)(E)=E'$).
% slanted
\margin{\uline{NB} $\widetilde\pi\backslash B$ est la cat. cofibrée sur $B$ définie par le foncteur composé $B\xrightarrow{\widetilde\pi}\mathrm{Ord}\longrightarrow\mathrm{Cat}$}
% ll. 10–12
Et \unsure{on a} un foncteur canonique
\[
 \begin{aligned}
 \widetilde\pi\backslash B&\longrightarrow(\underbrace{\pi\backslash B}_Z)^\wedge\\
 (b,\underset{\substack{\rotatebox{-90}{$\in$}\\\mathcal P^*(\pi(b))}}{E})&\longmapsto
 \widetilde E=\coprod_{a\in E}a\subset\coprod_{a\in\underset{\substack{\rotatebox{90}{$=$}\\\pi(b)}}{Z_b}}a .
 \end{aligned}\tag*{(*)}
\]
% ll. 13–19
Je voudrais des conditions sur $A$, ou encore sur $B=A^{\mathrm o}$, qui assurent que \inserted{quel que soit $\pi\in\operatorname{Ob}A^\wedge=\operatorname{Ob}B^\vee$,} le foncteur $(*)$ est pleinement fidèle$^*$, et son image essentielle est \deleted{\unsure{la plus|le plus}} stable par $\varprojlim$ finies. Par ailleurs, il est clair que cette image essentielle contient $Z$.

% ll. 20–21
Je m'aperçois qu'il est déraisonnable de s'attendre que le foncteur $(*)$ soit
% archive sheet number: 12

\scanpage{14}
\mspage{7}
% ll. 1–5
pleinement fidèle. Je vais par contre exiger que (pour tt $\pi$) la catégorie $\widetilde\pi\backslash \unsure{E|B}$ soit stable par $\varprojlim$ finies \inserted{non vides}, et que le foncteur $(*)$ soit exact à gauche.
% ll. 5–15
Prenant pour $\pi$ le foncteur constant de valeur $e$, on trouve que $\widetilde\pi\backslash B\simeq\pi\backslash B\simeq B$, on veut donc que $B$ soit stable par $\varprojlim$ finies \inserted{\uline{non vides}}, i.e. $A$ stable par $\unsure{\varinjlim|\underline{\lim}|\varprojlim}$ finies \uline{non vides}. Pour $B$, cela revient à la stabilité a) par produit binaire, et b) par produit fibré binaire (ou, au choix, par noyaux \unsure{des|de} doubles-flèches). Montrons qu'alors $\widetilde\pi\backslash B$ est stable par les \unsure{$\overline{\rm ms}$|mêmes} opérations.
% ll. 16–20
En fait, il faudra faire une hypothèse sur $\pi$, assurant que $\deleted{\unsure{f|c}}C=\widetilde\pi\backslash B$ est non seulement cofibré sur $B$, mais aussi \uline{fibré}. Cette hypothèse, c'est que $\pi$ transforme flèches de $B$ en
% archive sheet number: 13

\scanpage{15}
\mspage{8}
% letterhead
\begin{center}
ACADEMIE D'AIX-MARSEILLE\\
\rule{4em}{0.4pt}\\
INSPECTION ACADEMIQUE\\
DE VAUCLUSE\\
\rule{1.6em}{0.4pt}
\end{center}
% ll. 1–8
épimorphismes. Alors \deleted{l'\unsure{application|applicat.}} pour $u:b\longrightarrow b'$ donné, l'application
\[
\begin{aligned}
 u_*\deleted{\unsure{*|=}}=\widetilde\pi(u):\quad\widetilde\pi(b)&\longrightarrow\widetilde\pi(b'),\qquad E\longrightarrow\pi(u)(E),\\
 \rotatebox{90}{$=$}&\hspace{3em}\rotatebox{90}{$=$}\\
 \mathcal P^*(\pi(b))&\longrightarrow\mathcal P^*(\pi(b'))
\end{aligned}
\]
admet un adjoint à droite \deleted{$\unsure{\pi|u}$} $u^*$, savoir $E'\longrightarrow\pi(u)^{-1}(E')$, qui transforme un $E'$ \uline{non vide} en \deleted{$\unsure{F|I}$}$E$ non vide. (\unsure{because|beca-use} $\pi(b)\longrightarrow\pi(b')$ surjectif)
% ll. 9–15
\uline{Lemme 1} Soit $C\longrightarrow B$ une \uline{fibration}, dans $\mathrm{Cat}$, et \deleted{\unsure{suppos|suppos.}} soit $I$ dans $\mathrm{Cat}$ telle que les $\varprojlim$ de type $I$ existent dans $B$, et dans les fibres $C_b$. Alors elles existent dans $C$, et $C\longrightarrow B$ commute aux $\varprojlim$ de type $I$.
% ll. 16–19
\uline{Corollaire} Soit $\pi:B\longrightarrow\mathrm{Ens}$ un foncteur, \deleted{transformant flèches de $B$ en \unsure{épis|épi}}, soit
\diagram{\begin{tikzcd}[ampersand replacement=\&]
 \widetilde\pi:B\arrow[r,"\pi"]\&\mathrm{Ens}\arrow[r,"{\mathcal P^{\unsure{*|\epsilon}}}"]\&\mathrm{Ord}\arrow[d,"\text{can.}"]\\
 \&\&\mathrm{Cat}
\end{tikzcd}}
le composé, \deleted{\unsure{et|de}} Alors $C$
% archive sheet number: 14

\scanpage{16}
\mspage{9}
% ll. 1–6
catégorie cofibrée sur $B$, à fibres les $\widetilde\pi(b)=\deleted{\unsure{\mathcal P|P}}\mathrm{Cat}\,\mathcal P^{\deleted{\unsure{*|+}}}(\pi(b))$, définie par $\widetilde\pi$, est aussi fibrée sur $B$\deleted{\unsure{,|$_g$}}, \deleted{\unsure{si|et}} et elle est stable par tout type de $\varprojlim$ \deleted{\unsure{non vide|finies}} \inserted{(\unsure{proj. ind.|prop. ind.|resp. ind.})} $I$ \deleted{tel que} pour lequel $B$ \unsure{lui-\,$=$|lui-même} est stable.
% archive sheet number: 15

\scanpage{18}
\mspage{0}
% ll. 1–3
$\deleted{\unsure{P|\Gamma}}\Delta(E)={}$ groupe engendré par les éléments de la forme $[x,y]$ $(x,y)\in E$, avec les relations
\[
[x,x]=1,\quad [x,y][y,x]=1,\quad [x,y][y,z]=[x,z].
\]

% ll. 4–6
NB Je soupçonne que $\Delta(E)$ est isomorphe au sous-groupe du groupe libre $L(E)$ engendré par $E$, engendré par les éléments de la forme $[x,y]=x^{-1}y$.

% ll. 7–13
Si $\operatorname{card}E=n$ est fini, $\Delta(E)$ est un groupe libre à $n-1$ générateurs, si $E=[1,n]$, il est engendré par les générateurs libres $\deleted{\unsure{E_3|E_{\supset}|E_i}}$ $g_i=[i,i+1]$ pour $i\in[1,n-1]$. On a alors pour $i<j$
% chain at left, beside lines 8–10
\diagram{\begin{tikzcd}[ampersand replacement=\&,column sep=small]
1\arrow[r,dash]\&2\arrow[r,dash]\&3\arrow[r,dotted,dash]\&n-1\arrow[r,dash]\&n
\end{tikzcd}}
\[
\left\{\begin{aligned}
[i,j]&=g_ig_{i+1}\ldots g_{j-1}\\
[j,i]&=[i,j]^{-1}=g_{j-1}^{-1}\deleted{\unsure{g|j}}\ldots g_{i+1}^{-1}g_i^{-1}
\end{aligned}\right.
\]

% ll. 14–16
Résultat analogue si $E$ est dénombrable, alors $\Delta(E)$ est un groupe libre avec un ens dénombrable de générateurs libres.

% ll. 17–24
Soit $X$ une catégorie ayant un plus petit objet $e$, et un plus grand objet $a$. Alors $\pi_1(X,e)$ est isomorphe \deleted{\unsure{au|à un} \unsure{groupe|gr}} à un quotient du groupe $\Delta(\operatorname{Hom}(e,a))$, par les relations suivantes : pour tout système de doubles flèches
\diagram{\begin{tikzcd}[ampersand replacement=\&]
e\arrow[r,shift left=.6ex,"u"]\arrow[r,shift right=.6ex,"v"']\&x\arrow[r,shift left=.6ex,"\alpha"]\arrow[r,shift right=.6ex,"\beta"']\&a
\end{tikzcd}}
, \deleted{on} on a la relation
\[
\boxed{[\alpha v,\alpha u]=[\beta v,\beta u]}\quad\text{i.e.}\quad
[\alpha v,\alpha u][\beta v,\beta u]^{-1}=1.
\]
% archive sheet number: 17

\scanpage{19}
% Three long oblique cancellation strokes run across the sheet.
% l. 1
\[
\underline{x}=(x_i)_{i\in I}\qquad \deleted{{\unsure{Y|y}}^{\wedge}}
\]

% ll. 2–5, left
$\deleted{\unsure{\operatorname{Hom}|h}}\quad \deleted{\unsure{a|b}}\quad
\deleted{\unsure{F|E}}\qquad X_{//a}\qquad X_{//\underline{a_i}}\qquad X_{//a}$

% ll. 2–4, right
\diagram{\deleted{\hbox{\begin{tikzcd}[ampersand replacement=\&]
\&z\arrow[r]\&x\arrow[r,shift left=.6ex]\arrow[r,shift right=.6ex]\&\unsure{y|y'}\\
z'\arrow[ur]\arrow[r]\&x'\arrow[ur]\&\&
\end{tikzcd}}}}

% ll. 5–8
\diagram{\deleted{\hbox{\begin{tikzcd}[ampersand replacement=\&]
{}\arrow[r,Rightarrow]\&x_1\arrow[r,shift left=.55ex]\arrow[r,shift right=.55ex]\&x_0
\end{tikzcd}}}}
\diagram{\deleted{\hbox{\begin{tikzcd}[ampersand replacement=\&]
\&a_i\arrow[dr,dash]\&\\
\unsure{g|b|a}\arrow[ur]\arrow[rr]\arrow[dr]\&\&a_{\unsure{ij|j}}\\
\&a_j\arrow[ur]\&
\end{tikzcd}}}}
\diagram{\begin{tikzcd}[ampersand replacement=\&]
\substack{x\\a}\arrow[r,shift left=.5ex,"f"]\arrow[r,shift right=.5ex,"g"']\&\deleted{\unsure{y|g}}
\end{tikzcd}}
$\deleted{\unsure{g|q}}$
\diagram{\begin{tikzcd}[ampersand replacement=\&]
b\arrow[r,shift left=.5ex,"f"]\arrow[r,shift right=.5ex,"g"']\&{}
\end{tikzcd}}
\diagram{\deleted{\hbox{\begin{tikzcd}[ampersand replacement=\&]
\&\&x\arrow[dr]\&\deleted{\unsure{a|x}}\\
\&y\arrow[ur,shift left=.5ex]\arrow[ur,shift right=.5ex]\arrow[rr]\&\&a\\
z\arrow[ur]\arrow[urrr]\&\&\&
\end{tikzcd}}}}

% ll. 9–11
\diagram{\deleted{\hbox{\begin{tikzcd}[ampersand replacement=\&,column sep=large]
{}\arrow[r,shift left=.55ex,"{\deleted{\unsure{a|a_1}}}"]\arrow[r,shift right=.55ex]\&
{}\arrow[r,shift left=.55ex,"u'"]\arrow[r,shift right=.55ex,"v'"']\&
{}\arrow[r,shift left=.55ex,"u"]\arrow[r,shift right=.55ex,"v"']\&{}
\end{tikzcd}}}}
$\deleted{\unsure{uu',vv'|uu'=vv'}}\quad vu',uv'\qquad \deleted{\unsure{H|M}}$
\[
H_{m,n}\times H_{n,p}\longrightarrow H_{m,p}
\]
\diagram{\begin{tikzcd}[ampersand replacement=\&]
x\arrow[d]\\ H_x=\unsure{a|x}/R_x
\end{tikzcd}}

% ll. 12–15
\[
\begin{aligned}
fu&=gv\\
gv'&=gw\\
fu v_1&=g\underbrace{v v_1}=g\underbrace{v'v'_1}=gwv'_1
\end{aligned}
\]
\diagram{\deleted{\hbox{\begin{tikzcd}[ampersand replacement=\&]
z\arrow[r]\&x\arrow[r,shift left=1ex,"f"]\arrow[r,"g" description]\arrow[r,shift right=1ex,"h"']\&y\\
z'\arrow[ur,shift left=.3ex]\arrow[ur,shift right=.3ex]\&\&
\end{tikzcd}}}}
\diagram{\begin{tikzcd}[ampersand replacement=\&]
x\arrow[r]\arrow[d]\&y\arrow[d]\\ H_x\arrow[r,dash]\arrow[d,"{\mathrm{iso}\,?}"]\&H_y\\
[x]\&
\end{tikzcd}}
% archive sheet number: 18
